\section{Conclusion}

The mathematics of a listed derivatives market is organized by the objects its rules transform. An allocation map consumes resting quantity and assigns fills. A leg map converts strategy trades into signed positions. A resource matrix prevents double use of orders supporting implied opportunities. An observation map determines which aspects of those mechanisms are visible. Opening and settlement rules select prices for different purposes. A monetary-value map translates positions and trades into exact variation. Margin and collateral maps determine whether the resulting positions can be maintained. Exercise and delivery create or discharge additional obligations.

Several conclusions follow from keeping those maps separate. Integer pro-rata allocation need not compose across events. Implied quantities need not add. An MBO sequence need not imply FIFO execution. A futures quote-value integral is not an upfront purchase payment. Terminal profit does not determine peak funding. Convex portfolio risk can rise when one hedge leg is removed. Enough aggregate collateral value does not guarantee enough settlement cash. A final futures mark and a physical invoice are different ledger entries.

The framework's mathematical building blocks are deliberately transparent. Prefix geometry establishes FIFO and waterfall conservation. Convexity describes depth costs and the declared scenario-risk model. Linear constraints and dual certificates expose shared-resource bottlenecks. Telescoping identifies exact settlement totals; prefix maxima identify their funding requirements. None of these tools removes the need to specify the contract, market phase, account, and operative rule branch.

Further work can add empirical event intensities, full product configurations, calibrated option repricing, a complete implied-recipe priority engine, and time-dependent collateral transformation. Such extensions should preserve the distinctions and reconciliations established here. A more realistic model is useful only if its additional detail enters the correct part of the lifecycle.
